% Combined submission report; Task 1 response text retained from the existing draft.
% Compile from the repository root so the numbered figure PDFs can be found.
% The included Task 1 notebook has completed full-preset outputs and implemented exercises.
% Exported CSVs support the reported numbers.
\documentclass[11pt,a4paper]{article}
\usepackage[margin=22mm]{geometry}
\usepackage[T1]{fontenc}
\usepackage{lmodern}
\usepackage{microtype}
\usepackage{amsmath,amssymb}
\usepackage{booktabs,tabularx}
\usepackage{graphicx}
\usepackage{caption}
\usepackage[hidelinks]{hyperref}
\usepackage{enumitem}
\setlength{\parindent}{0pt}
\setlength{\parskip}{6pt}
\setlist{nosep,leftmargin=*}
\captionsetup{font=small,labelfont=bf}
\newcommand{\studentname}{Burhan Ahmad}
\newcommand{\studentid}{27100405}
\newcommand{\repository}{\url{https://github.com/BuriAhmad/AI651-Assignment1}}
\newcommand{\rmseunit}{\ensuremath{\mathrm{m/s^2}}}
\newcommand{\mseunit}{\ensuremath{(\mathrm{m/s^2})^2}}
\newcommand{\evidencefigure}[3]{%
  \begin{figure}[htbp]
    \centering
    \includegraphics[width=\linewidth,height=.78\textheight,keepaspectratio]{#1}
    \caption{#2}\label{#3}
  \end{figure}}

\begin{document}
\begin{center}
  {\Large\bfseries AI651 --- Assignment 1}\par
  {\large Tasks 1 and 2}\par
  \studentname\quad |\quad \studentid
\end{center}
\textbf{Code repository:} \repository

\section*{Task 1: Forecasting Across Heterogeneous Sensors}
\textbf{Experimental setup.} A 96-sample context from one vibration sensor is used to forecast 48 samples at 20 Hz. Stations 1--3 are established; Station 4 is excluded from model fitting and scale estimation but supplies its own context at inference. Training, validation, and final-test target regions are $[0,9600)$, $[9600,12480)$, and $[12480,15360)$, respectively. The results below use the supplied full-run exports. Unless explicitly labeled otherwise, errors are on validation data; Task 1 RMSEs are in \rmseunit{} and MSEs in \mseunit{}.

\section*{Response 1A}
Before Output 1.1, my predicted ordering for established stations was Period-routed ridge $>$ Raw Attention $>$ Sensor-specific ridge $>$ Shared ridge. For held-out Station 4, I predicted Period-routed ridge $>$ Raw Attention $>$ Shared ridge, with Sensor-specific ridge unavailable.

The measured established-station ordering was Period-routed ridge (RMSE 0.166) $>$ Raw Attention (0.424) $>$ Shared ridge (0.767) $>$ Sensor-specific ridge (0.768). I therefore revise only the two fixed ridges: their errors are almost tied, and the tiny difference does not establish a meaningful advantage. Period-routed ridge was best for every product in both populations (Table~\ref{tab:products}). Sensor-specific ridge was slightly better than Shared ridge on established Products A and B, but worse on C. On held-out Station 4, my predicted ordering was confirmed: Period-routed ridge (0.173) $>$ Raw Attention (0.612) $>$ Shared ridge (0.876).

Period routing reduced RMSE relative to Raw Attention by approximately 61\% on established stations and 72\% on Station 4. This supports operating pace as the main conditioning variable. Raw Attention builds its mixing from the current context, while Shared ridge and each Sensor-specific ridge matrix apply fixed operators that must compromise across different paces. Conditioning only on station identity does not resolve changes of pace within a station.

\begin{table}[htbp]
\centering\small
\caption{Product-level validation RMSE from Output 1.1. Established means Stations 1--3; unavailable entries reflect the absence of a Sensor-specific ridge fit for Station 4.}
\label{tab:products}
\begin{tabular}{lrrr rrr}
\toprule
&\multicolumn{3}{c}{Established}&\multicolumn{3}{c}{Held-out Station 4}\\
\cmidrule(lr){2-4}\cmidrule(lr){5-7}
Model&A&B&C&A&B&C\\
\midrule
Shared ridge&0.805&0.854&0.621&0.909&0.982&0.713\\
Sensor-specific ridge&0.804&0.852&0.630&---&---&---\\
Period-routed ridge&0.181&0.164&0.152&0.184&0.182&0.150\\
Raw Attention&0.385&0.461&0.422&0.645&0.562&0.627\\
\bottomrule
\end{tabular}
\end{table}

\section*{Response 1B}
Output 1.2 shows that station identity alone does not determine the continuation rule. Station 1 has a period of 16.135 samples for Product A and 30.093 for Product C. One fixed ridge matrix must compromise between these different oscillation rates. In the illustrated validation windows, the fixed ridge forecasts fail to track the peaks and troughs of both regimes as closely as Period-routed ridge (Figure~\ref{fig:forecasts}); its window RMSEs are 0.232 and 0.150 for A and C, respectively.

Attention's learned projections $W_Q,W_K,W_V$ are shared, but $Q=XW_Q$, $K=XW_K$, and $V=XW_V$ depend on the context $X$. Changing $X$ changes the query--key comparisons, the Attention weights, and therefore the value-mixing operation. Output 1.3 confirms this: the two contexts have mean absolute weight difference 0.011 and maximum difference 0.410 (Figure~\ref{fig:attention}). Ridge instead applies its same fitted operator to every context assigned to that matrix. These maps show context-dependent mixing; they do not establish recovery of the physical operating period.

\section*{Response 2}
The selected moving-average width was 41 samples. Output 2.1 shows how candidate widths change the trend and remainder: the wider average leaves the faster oscillations more visible in the remainder, and width 41 produces a sharp period-score peak near the actual period (Figure~\ref{fig:decomposition}). Output 2.2 quantifies this on established-station \emph{training} windows: within-one-sample recovery rises from 0.687 for raw signals to 0.998 at width 41. Width 41 maximizes this recovery measure, although width 25 has the smaller median absolute period error (0.122 versus 0.216 samples).

The oracle diagnostic is useful because the synthetic generator provides the true period for analysis, allowing width selection against known repetition using training data only. This truth is not supplied to the forecasting model. Output 2.3 also shows a forecasting improvement: established validation RMSE/MSE falls from 0.424/0.180 for Raw Attention to 0.298/0.089 for Attention + decomposition. Held-out validation RMSE/MSE falls from 0.612/0.375 to 0.445/0.198.

A moving average is a sensible bias for a slowly varying baseline superimposed on faster vibration. The remainder is an estimate, not the true generating component. Abrupt baseline changes would violate the smoothness assumption: smoothing would smear the discontinuity over neighboring positions and contaminate the remainder.

\section*{Response 3}
Output 3.1 gives FFT-based delay scores $4,-4,4,-4$ for delays 0--3, matching direct circular correlation. The required aggregation example uses values $[10,20,30,40]$, delays $[1,3]$, and weights $[0.75,0.25]$. With
\[
 z_t=\sum_j\alpha_jv_{(t-\tau_j)\bmod L},
\]
the expected output is $[35,15,25,25]$; in particular, $z_0=0.75(40)+0.25(20)=35$, establishing the $t-\tau$ indexing direction.
% The saved notebook does not contain executed aggregation-check evidence.
% Replace "expected output" with "observed output" only after checking the executed example.

In Output 3.2, the illustrated Station 2 Product A run has $P=16.135$ samples. Strong residual recurrence occurs at delays 16 and 32, with correlations 0.896 and 0.897, close to $P$ and $2P=32.269$ (Figure~\ref{fig:recurrence}). Slow drift makes the raw signal similar to shifted copies over many lags, obscuring the distinctive recurring delays. Removing it exposes sharper peaks and supports the delay mixer's displacement-based bias.

This assumption could fail for an isolated transient or a period changing rapidly within one context: no small set of recurring displacements would remain informative, and shifted values could combine unrelated events. Output 3.2 is a signal-level correlation diagnostic, not the trained mixer's learned query/key scores.

\section*{Response 4}
\textbf{(a)} Before Output 4.1, I predicted that decomposition would help most when the slow component was prominent because it separates slow trend from repetition before forecasting.

\textbf{(b)} Established validation errors (Outputs 1.1 and 4.1):
\begin{center}
\small
\begin{tabular}{lrr}
\toprule
Model & RMSE & MSE \\
\midrule
Shared ridge & 0.767 & 0.588 \\
Sensor-specific ridge & 0.768 & 0.590 \\
Period-routed ridge & 0.166 & 0.028 \\
Raw Attention & 0.424 & 0.180 \\
Attention + decomposition & 0.298 & 0.089 \\
Raw delay mixer & 0.223 & 0.050 \\
Autoformer-inspired & 0.199 & 0.040 \\
\bottomrule
\end{tabular}
\end{center}

Decomposition improved Attention (0.424 to 0.298) and delay mixing (0.223 to 0.199); delay mixing improved raw (0.424 to 0.223) and decomposed inputs (0.298 to 0.199). Established RMSE gains overlapped rather than adding. This is an empirical observation, not a causal explanation. Errors generally increased across the horizon; Period-routed ridge stayed lowest, while Autoformer-inspired was strongest among neural models overall (Figure~\ref{fig:horizon}).

\textbf{(c)} I selected Period-routed ridge for both populations: validation RMSEs were 0.166/0.173, with 59,904 parameters and 0.003\,s fitting time. Autoformer-inspired had worse validation RMSEs (0.199/0.241), 373,026 parameters, and 819.367\,s fitting time. Output 4.3 gave established/held-out test RMSEs of 0.161/0.183 and MSEs of 0.026/0.033, supporting the validation choices.

\clearpage
\section*{Task 2: Autoformer forecasting challenge}
\subsection*{Problem, data, and modeling evidence}
The observed target contains 43,656 regularly spaced samples. The final forecast covers exactly indices 43,657--43,824, a 168-step horizon. The ten aligned external measurements contain six continuous variables (A--F) and four mutually exclusive indicators (G--J). Both historical and future-horizon external values are supplied; future target values remain hidden. Structural checks confirm index continuity, nonmissing nonnegative training targets, hidden test targets, full external coverage, and one-hot indicators. Task 2 errors are in the anonymized target's original units, not Task 1's vibration units; sMAPE is a percentage.

Diagnostics present in the notebook were checked against the supplied data. The target ranges from 0 to 994, with mean 98.17, median 73, standard deviation 90.83, skewness 1.83, and excess kurtosis 5.00. The 99th percentile is 418. The correlation between 168-step rolling mean and standard deviation is 0.787 for raw values and $-0.180$ after $\log(1+y)$, supporting level-dependent variance and a scale-stabilizing transform. This motivates testing a log transform; it does not establish that it will optimize original-scale RMSE.

Global ACF is 0.966 at lag 1, 0.401 at lag 24, 0.158 at lag 48, and 0.028 at lag 168. Local lag-24 ACF over 672-point windows varies from $-0.026$ to 0.637. The periodogram and local ACF were inspected rather than treating one spectral peak as a fixed operating period. The evidence supports strong short-lived dependence and changing local recurrence (Figure~\ref{fig:t2acf}), so copying Task 1's fixed pace routing would lack the same justification.

External associations are modest: contemporaneous correlations are 0.169 for A and $-0.245$ for D; they are diagnostic associations rather than causal claims. The supplied future regime profile differs from history: H occurs in 51.79\% of future steps versus 32.21\% historically, while I falls to 8.93\% from 34.99\%. This supports testing known-future covariates and also cautions against assuming the final horizon has the historical average regime mix.

\subsection*{Chronological validation and preprocessing}
Screening uses seed 2026 and origins 43,320 and 43,488, at most four epochs, patience one, and training-window stride eight. Final evaluation uses seeds 2026 and 3407 and origins 43,152, 43,320, and 43,488. Their nonoverlapping targets are 43,153--43,320, 43,321--43,488, and 43,489--43,656. Each origin permits target history only through that origin. The final stage uses stride four, at most ten epochs, patience two, and 10,663/10,705/10,747 training windows at the three origins. Shuffling eligible training windows does not randomize the temporal holdout.

Each fold fits target location/scale and continuous-covariate scalers solely on pre-origin history. Indicators remain binary. Known-future covariates enter the decoder for the validation horizon; future target values never enter it. Metrics are computed after inversion to original target units and clipping forecasts to nonnegative values. Two seeds and three folds expose initialization and temporal variation at manageable cost; these six evaluations are not six independent datasets.

\subsection*{Autoformer implementation and selection}
The self-contained encoder--decoder follows the decomposition and autocorrelation ideas of Autoformer~\cite{autoformer}, extending the mechanisms explored in Task 1. Moving averages with replicated endpoints split hidden sequences repeatedly after encoder mixing/feed-forward sublayers and decoder self/cross mixing/feed-forward sublayers. Decoder trend contributions are projected and accumulated rather than discarded. Seasonal layer normalization removes the temporal mean.

The mixer takes FFTs of learned queries and keys, forms circular delay correlations, averages scores across heads/features, selects approximately $3\ln L$ delays per example, and softmax-weights shifted values. It therefore aggregates temporal displacements rather than constructing vanilla pairwise dot-product Attention. Encoder inputs contain target history plus projected historical covariates. The decoder starts with 84 historical seasonal/trend positions, zero future seasonal targets, and a history-mean future trend; it also receives the known future covariates. No future target or inferred calendar identity is used.

The question-driven search first tests target representation, then history/decomposition, then capacity; all screening comparisons use the same seed/folds (Table~\ref{tab:t2screen}). Raw global standardization wins mean RMSE 50.49 versus 61.58 for log1p, although log1p has lower sMAPE (55.05 versus 83.49). This demonstrates metric disagreement: stabilizing relative scale need not improve squared error on large target levels. Context 336/kernel 25 wins the structure stage; extending to 672/kernel 49 does not help these folds. Width 96 then improves screening RMSE from 50.49 to 48.36 while increasing parameters from 117,889 to 262,849. These are screening observations, not multiple-seed architecture superiority claims.

\begin{table}[htbp]
\centering\small
\caption{Screening means over one seed and two chronological horizons. Candidate names encode context length, kernel, or model width.}\label{tab:t2screen}
\begin{tabular}{llrrrr}
\toprule
Stage & Candidate & RMSE & MAE & sMAPE (\%) & Parameters \\
\midrule
Target & T\_raw & 50.49 & 36.40 & 83.49 & 117,889 \\
Target & T\_log1p & 61.58 & 36.07 & 55.05 & 117,889 \\
Structure & S\_168\_k25 & 54.37 & 38.35 & 94.66 & 117,889 \\
Structure & S\_336\_k25 & 50.49 & 36.40 & 83.49 & 117,889 \\
Structure & S\_336\_k49 & 52.43 & 36.51 & 96.10 & 117,889 \\
Structure & S\_672\_k49 & 55.98 & 43.72 & 86.09 & 117,889 \\
Capacity & C\_compact32 & 52.76 & 36.93 & 79.64 & 30,273 \\
Capacity & C\_base64 & 50.49 & 36.40 & 83.49 & 117,889 \\
Capacity & C\_medium96 & 48.36 & 33.51 & 84.22 & 262,849 \\
\bottomrule
\end{tabular}
\end{table}

Selection minimizes mean fold RMSE. Candidates within 0.5\% of the best are ordered by parameter count and then mean best epoch before RMSE. The unknown leaderboard penalty coefficients are not estimated, and no leaderboard score is used as a validation signal.

\subsection*{Required external-data ablation and final results}
The final comparison holds the selected architecture and optimization settings fixed, changing only external inputs. Both variants use context 336, label length 84, horizon 168, model width 96, four heads, two encoder layers, one decoder layer, feed-forward width 192, moving average 25, factor 3, and dropout 0.1. AdamW uses learning rate and weight decay $10^{-4}$, batch size 32, standardized raw-target MSE loss, and gradient clipping at 1.0.

\begin{table}[htbp]
\centering\small
\caption{Final ablation: arithmetic means over two seeds $\times$ three folds; RMSE spread is the sample standard deviation across the six run-level RMSEs.}\label{tab:t2ablation}
\begin{tabular}{lrrrrr}
\toprule
External inputs & Mean RMSE $\pm$ SD & MAE & sMAPE (\%) & Parameters & Min/run\\
\midrule
A--J & $58.72\pm16.23$ & 45.25 & 82.61 & 262,849 & 27.65\\
None & $88.41\pm23.55$ & 76.54 & 102.41 & 260,929 & 25.96\\
\bottomrule
\end{tabular}
\end{table}

External inputs reduce mean RMSE by 29.69 (33.58\%), MAE by 31.29 (40.89\%), and sMAPE by 19.81 percentage points. Their extra 1,920 parameters are a 0.74\% increase. Every matched fold/seed RMSE improves (Table~\ref{tab:t2folds}), supporting selection of \texttt{FINAL\_with\_external} under the notebook's rule. This is the measured benefit of all A--J jointly; the experiment does not identify the individual contribution of each feature or separate historical from future-covariate benefit.

\begin{table}[htbp]
\centering\small
\caption{All six final fold/seed pairs. Epoch is the selected checkpoint epoch with external data; the complete per-variant metrics and epoch fields are retained in \texttt{final\_cv\_results.csv}.}\label{tab:t2folds}
\begin{tabular}{rrrrrrr}
\toprule
Origin & Seed & Epoch & \shortstack{RMSE\\with} & \shortstack{MAE\\with} & \shortstack{sMAPE (\%)\\with} & \shortstack{RMSE\\without}\\
\midrule
43152 & 2026 & 4 & 71.11 & 54.76 & 74.26 & 117.06 \\
43152 & 3407 & 4 & 77.09 & 63.29 & 78.22 & 119.52 \\
43320 & 2026 & 2 & 60.27 & 43.98 & 81.40 & 76.25 \\
43320 & 3407 & 1 & 65.73 & 54.44 & 94.92 & 79.93 \\
43488 & 2026 & 2 & 38.48 & 29.47 & 89.28 & 68.64 \\
43488 & 3407 & 5 & 39.62 & 25.54 & 77.56 & 69.04 \\
\bottomrule
\end{tabular}
\end{table}

Temporal difficulty dominates much of the spread. With-external fold-mean RMSEs are 74.10, 63.00, and 39.05 in chronological order; seed-mean RMSEs are 56.62 and 60.81. This favors reporting all folds rather than presenting the easiest final horizon or one seed as overall performance.

Horizon curves pool the six squared errors at each step, then take a square root (Figure~\ref{fig:t2horizon}); the pooled all-step RMSE is 60.56 with external inputs versus 90.98 without, distinct from the arithmetic mean run RMSE in Table~\ref{tab:t2ablation}. With-external block RMSE rises from 38.64 over steps 1--24 to 64.01 over steps 145--168 and reaches 74.97 at steps 121--144. Neither curve increases monotonically, and the external advantage is not uniform at every step.

\subsection*{Full-history training, predictions, and limits}
The final model is freshly initialized with seed 2026 and fitted on all 43,656 observed targets. Its training-window stride is two, with scalers fitted on the full observed history. The six with-external best epochs are 4, 2, 2, 4, 1, and 5; their median selects \textbf{three full-refit epochs}. The checkpoint confirms \textbf{262,849 trainable parameters}; refitting took 34.71 minutes. The saved prediction array and CSV contain exactly 168 finite nonnegative values in chronological order. They are preserved without regeneration; 24 values are zero after clipping. The original forecast is shown in Figure~\ref{fig:t2forecast}. Hidden-target performance is unavailable, and no leaderboard submission was made during preparation.

The epoch ledger distinguishes the final three epochs from preceding training. The original early-stop loops imply 55 screening epochs and 58 final-CV epochs (30 with and 28 without external data), plus one smoke epoch and three refit epochs: \textbf{117 epochs for the complete pipeline}. The prepared metadata uses this conservative whole-pipeline count rather than silently treating three refit epochs as the entire cost. Counts are derived from saved best epochs, patience, and epoch caps; the ledger states that derivation explicitly.

The original notebook does not call \texttt{model.eval()} in validation or final prediction. Consequently, dropout remains active even under \texttt{no\_grad}; reported metrics and forecasts describe that stochastic inference protocol. Also, the same fold targets determine checkpoints and compare architectures, so these are model-selection estimates rather than an independent final test. The three recent folds, changing local recurrence, and shifted future covariate mix limit generalization claims. No expensive rerun or silent replacement of results was performed. The observed external-data advantage supports this selected run, while deterministic evaluation and untouched additional origins would strengthen a future confirmation.

As in Task 1, useful model structure must match the data. Here, local decomposition/delay mixing is paired with legitimately known future measurements, while transform, context, and capacity are chosen by chronological evidence rather than copied from the synthetic benchmark.

\evidencefigure{Question 2 - Leaderboard/figures/global_acf.png}{Task 2: original Kaggle global ACF diagnostic. Short-lag dependence decays substantially before the full 168-step horizon.}{fig:t2acf}
\evidencefigure{Question 2 - Leaderboard/figures/horizon_comparison.pdf}{Task 2: horizon RMSE derived from the saved \texttt{final\_horizon\_rmse.csv}, using a shared vertical scale for both final ablation variants.}{fig:t2horizon}
\evidencefigure{Question 2 - Leaderboard/figures/final_forecast.png}{Task 2: original Kaggle full-history forecast. The continuation has no revealed ground truth; it is not a test-accuracy plot.}{fig:t2forecast}


\clearpage
\appendix
\section*{Task 1 supporting numbered figures}
The following figures are the original exports used to interpret the responses. They retain their notebook Output identifiers. Output 4.2 summarizes the full 48-step horizon; its curves fluctuate, so the response describes the overall increase rather than claiming monotonic growth.

\evidencefigure{Question 1/results/design/1.2-1.pdf}{Output 1.2: Station 1 validation forecast cases under Products A and C. The same station exhibits different operating periods.}{fig:forecasts}
\evidencefigure{Question 1/results/design/1.3-1.pdf}{Output 1.3: Attention maps for two contexts and their difference. Shared learned projections produce different input-dependent weights.}{fig:attention}
\evidencefigure{Question 1/results/design/2.1-1.pdf}{Output 2.1: candidate moving-average widths, their residuals, and period-score evidence. Width 41 is selected.}{fig:decomposition}
\evidencefigure{Question 1/results/design/3.2-1.pdf}{Output 3.2: raw and residual signal correlations for a Product A validation context. The peaks are diagnostic recurrence evidence, not learned query/key scores.}{fig:recurrence}
\evidencefigure{Question 1/results/design/4.2-1.pdf}{Output 4.2: validation horizon RMSE for established stations and held-out Station 4.}{fig:horizon}

\clearpage
\section*{AI-use record for Task 1 report preparation}
Generative AI assisted with interpretation and drafting of the responses supplied in chat. Codex checked the reported numbers against the exported tables and formatted the material into this LaTeX document.

\textbf{Prompt excerpt supplied to Codex:} ``then i want you to create the latex and write resposes inside it. all in one go okay?'' This was followed by the Task 1 interpretation and the five drafted responses.

\textbf{Output:} this Task 1 LaTeX draft, containing Responses 1A, 1B, 2, 3, and 4, a product-level results table, and numbered supporting figures.

\textbf{Edits in this preparation step:} numerical claims were compared with CSV exports; held-out product results and the training-only width-selection qualification were added; Response 4 was shortened and its switch comparisons made explicit; the unverified assertion that the aggregation check passed was changed to the expected numerical example.

\textbf{Disclosure limit:} the earlier prompts that produced the supplied interpretation draft were not included in the material available for report preparation. This record describes the available interaction without reconstructing unknown prompts.

\section*{AI-use record for Task 2 preparation}
\textbf{Prompt:} ``Finish AI651 Assignment 1 --- Task 2 from the completed Kaggle run, update the report, and prepare the submission repo.'' The detailed instruction specified downloading the saved run, using all fold/seed artifacts, preserving Task 1, excluding optional reflection questions, and publishing a cleaned public repository.

\textbf{Output:} Codex recovered the saved version 2 bundle, checked numeric summaries and the checkpoint metadata, generated an epoch ledger and a shared-axis plot from existing results, prepared the 168-value paste list, and integrated Task 2 into this report.

\textbf{Edits and checks:} the briefing's example (260,929 parameters and RMSE 69.04) was identified as a without-external row rather than overall winning performance. Reporting uses the actual 262,849-parameter winner and six-run aggregates. The active-dropout limitation and shared validation/selection horizons are disclosed. The original Task 2 notebook and forecasts were preserved, and no model was retrained or leaderboard feedback used.

\begin{thebibliography}{1}
\bibitem{autoformer} Wu et al. (2021). \emph{Autoformer: Decomposition Transformers with Auto-Correlation for Long-Term Series Forecasting}, Sections 3.1--3.2. Architectural reference listed in the assignment handout; the submitted notebook is self-contained and builds on the Task 1 mechanisms.
\end{thebibliography}

\end{document}
